\chapter*{LỜI NÓI ĐẦU}
\addcontentsline{toc}{chapter}{LỜI NÓI ĐẦU}

Thực tập kỹ thuật là học phần giúp sinh viên kết nối kiến thức cơ sở ngành với
môi trường làm việc thực tế. Đối với ngành Kỹ thuật Điện tử Viễn thông, việc
tiếp cận linh kiện điện tử, thiết bị tự động hóa, tài liệu kỹ thuật và quy trình
lựa chọn giải pháp có ý nghĩa quan trọng trong quá trình hình thành tư duy kỹ
sư. Đây cũng là cơ hội để sinh viên rèn luyện tác phong làm việc, kỹ năng tìm
kiếm thông tin, đọc tài liệu chuyên ngành và trình bày kết quả kỹ thuật.

Trong học kỳ \Semester, em thực hiện học phần Thực tập Kỹ thuật tại
\CompanyName, dưới sự hướng dẫn tại đơn vị của anh \CompanySupervisor\ và sự
hướng dẫn của giảng viên \UniversitySupervisor. Báo cáo tập trung giới thiệu
khái quát doanh nghiệp, các lĩnh vực hoạt động có liên quan đến chuyên ngành,
đồng thời hệ thống hóa nội dung tìm hiểu về linh kiện điện tử và thiết bị tự
động hóa.

Thuận lợi lớn nhất là \CompanyName\ có danh mục linh kiện điện tử và thiết bị
tự động hóa khá đa dạng, kèm trang giới thiệu và trang bán hàng công khai. Nhờ
đó em có thể đối chiếu mô tả sản phẩm với datasheet của hãng, đồng thời được
anh \CompanySupervisor\ và giảng viên \UniversitySupervisor\ định hướng phạm vi
tìm hiểu cho sát chuyên ngành. Khó khăn chủ yếu đến từ thời gian thực tập
ngắn, danh mục mã hàng lớn và cách đặt mã khác nhau giữa các hãng. Em cũng
không tiếp cận được quy trình nội bộ, dữ liệu kho hay phòng thí nghiệm để đo
kiểm, nên phần phân tích ứng dụng trong báo cáo mới dừng ở mức phương án trên
giấy. Việc đọc tài liệu tiếng Anh và không bỏ sót hậu tố, kiểu vỏ, dải nhiệt độ
hay trạng thái vòng đời cũng đòi hỏi em phải tra cứu nhiều lần mới quen.


Em xin chân thành cảm ơn Ban lãnh đạo \CompanyName, anh \CompanySupervisor,
giảng viên \UniversitySupervisor\ cùng các cá nhân đã hỗ trợ em trong quá trình
thực hiện học phần.
